% Options for packages loaded elsewhere
\PassOptionsToPackage{unicode}{hyperref}
\PassOptionsToPackage{hyphens}{url}
\PassOptionsToPackage{dvipsnames,svgnames,x11names}{xcolor}
\documentclass[
  ngerman,
  10pt,
  a4paper,
]{article}
\usepackage{xcolor}
\usepackage[margin=22mm]{geometry}
\usepackage{amsmath,amssymb}
\setcounter{secnumdepth}{-\maxdimen} % remove section numbering
\usepackage{iftex}
\ifPDFTeX
  \usepackage[T1]{fontenc}
  \usepackage[utf8]{inputenc}
  \usepackage{textcomp} % provide euro and other symbols
\else % if luatex or xetex
  \usepackage{unicode-math} % this also loads fontspec
  \defaultfontfeatures{Scale=MatchLowercase}
  \defaultfontfeatures[\rmfamily]{Ligatures=TeX,Scale=1}
\fi
\usepackage{lmodern}
\ifPDFTeX\else
  % xetex/luatex font selection
  \setmainfont[Path=/Users/stefanhamann/.cache/codex-runtimes/codex-primary-runtime/dependencies/native/libreoffice-headless/libreoffice/LibreOfficeDev.app/Contents/Resources/fonts/truetype/,Extension=.ttf,BoldFont=DejaVuSans-Bold,ItalicFont=DejaVuSans-Oblique,BoldItalicFont=DejaVuSans-BoldOblique]{DejaVuSans}
  \setmonofont[Path=/Users/stefanhamann/.cache/codex-runtimes/codex-primary-runtime/dependencies/native/libreoffice-headless/libreoffice/LibreOfficeDev.app/Contents/Resources/fonts/truetype/,Extension=.ttf,BoldFont=DejaVuSansMono-Bold,ItalicFont=DejaVuSansMono-Oblique,BoldItalicFont=DejaVuSansMono-BoldOblique]{DejaVuSansMono}
\fi
% Use upquote if available, for straight quotes in verbatim environments
\IfFileExists{upquote.sty}{\usepackage{upquote}}{}
\IfFileExists{microtype.sty}{% use microtype if available
  \usepackage[]{microtype}
  \UseMicrotypeSet[protrusion]{basicmath} % disable protrusion for tt fonts
}{}
\makeatletter
\@ifundefined{KOMAClassName}{% if non-KOMA class
  \IfFileExists{parskip.sty}{%
    \usepackage{parskip}
  }{% else
    \setlength{\parindent}{0pt}
    \setlength{\parskip}{6pt plus 2pt minus 1pt}}
}{% if KOMA class
  \KOMAoptions{parskip=half}}
\makeatother
\ifLuaTeX
\usepackage[bidi=basic,shorthands=off]{babel}
\else
\usepackage[bidi=default,shorthands=off]{babel}
\fi
\ifPDFTeX
\else
\babelfont{rm}[Path=/Users/stefanhamann/.cache/codex-runtimes/codex-primary-runtime/dependencies/native/libreoffice-headless/libreoffice/LibreOfficeDev.app/Contents/Resources/fonts/truetype/,Extension=.ttf,BoldFont=DejaVuSans-Bold,ItalicFont=DejaVuSans-Oblique,BoldItalicFont=DejaVuSans-BoldOblique]{DejaVuSans}
\fi
\ifLuaTeX
  \usepackage{selnolig} % disable illegal ligatures
\fi
\setlength{\emergencystretch}{3em} % prevent overfull lines
\providecommand{\tightlist}{%
  \setlength{\itemsep}{0pt}\setlength{\parskip}{0pt}}
\usepackage{fvextra}
\RecustomVerbatimEnvironment{verbatim}{Verbatim}{breaklines=true,fontsize=\footnotesize}
\usepackage{xurl}
\usepackage{seqsplit}
\usepackage{adjustbox}
\usepackage{ragged2e}
\usepackage{titlesec}
\titleformat{\section}{\large\bfseries\raggedright}{}{0em}{}
\titleformat{\subsection}{\normalsize\bfseries\raggedright}{}{0em}{}
\DeclareRobustCommand{\codeword}[1]{\texttt{\seqsplit{#1}}}
\AtBeginEnvironment{longtable}{\small\RaggedRight}
\usepackage{newunicodechar}
\newunicodechar{𝒞}{\ensuremath{\mathcal C}}
\newunicodechar{𝔊}{\ensuremath{\mathfrak G}}
\newunicodechar{𝔤}{\ensuremath{\mathfrak g}}
\newunicodechar{𝓗}{\ensuremath{\mathcal H}}
\newunicodechar{ℋ}{\ensuremath{\mathcal H}}
\newunicodechar{𝔄}{\ensuremath{\mathfrak A}}
\newunicodechar{𝟙}{\ensuremath{\mathbf 1}}
\newunicodechar{𝕀}{\ensuremath{I}}
\newunicodechar{𝕋}{\ensuremath{\mathbb T}}
\newunicodechar{𝟘}{\ensuremath{\mathbf 0}}
\newunicodechar{⊞}{\ensuremath{\boxplus}}
\newunicodechar{∘}{\ensuremath{\circ}}
\newunicodechar{⟦}{\ensuremath{[\![}}
\newunicodechar{⟧}{\ensuremath{]\!]}}
\setlength{\emergencystretch}{4em}
\setlength{\parindent}{0pt}
\setlength{\parskip}{0.45em}
\AtBeginDocument{\hypersetup{colorlinks=true,linkcolor=black,urlcolor=blue}\pdfstringdefDisableCommands{\def\codeword#1{#1}}\RaggedRight}
\usepackage{fancyhdr}
\pagestyle{fancy}
\fancyhf{}
\fancyhead[L]{\small TFPT / Universalraum - v1.6.8}
\fancyfoot[R]{\thepage}
\setlength{\headheight}{15pt}
% The preserved full history now has three-digit page numbers.
\makeatletter
\renewcommand{\@pnumwidth}{2.6em}
\renewcommand{\@tocrmarg}{3.6em}
\makeatother
\usepackage{bookmark}
\IfFileExists{xurl.sty}{\usepackage{xurl}}{} % add URL line breaks if available
\urlstyle{same}
\hypersetup{
  pdflang={de-DE},
  colorlinks=true,
  linkcolor={Maroon},
  filecolor={Maroon},
  citecolor={Blue},
  urlcolor={Blue},
  pdfcreator={LaTeX via pandoc}}

\author{}
\date{}

\begin{document}

\section{TFPT / Universalraum einfach erklärt: dieselbe Maschine,
mehrere
Ansichten}\label{tfpt-universalraum-einfach-erkluxe4rt-dieselbe-maschine-mehrere-ansichten}

Stand v1.6.8, 15. September 2026.

\subsection{Wo wir stehen}\label{wo-wir-stehen}

Wir haben noch keine vollständige Theorie des Universums. Aber wir
können jetzt an einem kleinen Teil des ursprünglichen TFPT-Modells drei
Dinge genauer zusammenbringen: Was passiert darin? Was lässt sich davon
sehen? Und verändert ein gezielter Eingriff tatsächlich etwas anderes?

Stell dir eine kleine Maschine hinter zwei Fenstern vor. Das eine zeigt
Fermionpaare, das andere Bosonen. Beide sind hier zwei mögliche
Erscheinungsformen derselben erhaltenen Gesamtladung. Zwei
Momentaufnahmen können gleich aussehen, obwohl die Maschine innerlich
unterschiedlich schwingt. Dann sehen die Bilder etwas später verschieden
aus.

Der wichtige neue Schritt: Für einen genau bezeichneten Teil der
Maschine können wir die fehlende innere Schwingungsphase aus beiden
Ansichten zu drei passenden Zeiten zurückrechnen. Eine vierte Ansicht
wird daraus richtig vorhergesagt. Wir müssen dafür keine zusätzliche
große Maschine erfinden.

Das funktioniert mathematisch exakt unter klaren Voraussetzungen: Wir
kennen die Dynamik, können denselben Zustand wiederholt vorbereiten und
haben die benötigten Messungen. Gerade die Herkunft dieser Handgriffe
aus dem Compiler ist noch nicht vollständig gezeigt.

\subsection{Was die neu zugeschickte Arbeit
beiträgt}\label{was-die-neu-zugeschickte-arbeit-beitruxe4gt}

Sie schlägt vor, alle erlaubten Handlungsfolgen mitsamt ihren
Überlappungen als gemeinsame Beschreibung zu verwenden. Das ist ein
brauchbarer Ordnungsrahmen: Zwei vermeintliche Teillösungen müssen
dieselben Handlungen und Antworten wiedergeben, nicht bloß ähnliche
Zahlen enthalten.

Die mitgelieferten Rechnungen konnten wir wiederholen. Trotzdem gibt es
einen echten Textfehler: Die vollständige erste Zeile dieser
Überlappungstabelle enthält bereits die gesamte Information, wenn alle
zusammengesetzten Handlungsfolgen dazugehören. Man braucht nicht zwei
unabhängige riesige Datensammlungen. Nur einige wenige Tabellenwerte
reichen natürlich nicht.

Auch ihre Idee eines kleinen Ladungsspeichers trägt: Wenn ein Teil mehr
Ladung erhält, kann ein mitgerechneter Speicher entsprechend weniger
haben. Der Speicher darf danach aber nicht so behandelt werden, als wäre
nichts geschehen. Er ist Teil des Vorgangs und gegebenenfalls
verbraucht.

\subsection{Drei neue Ergebnisse in
Bildern}\label{drei-neue-ergebnisse-in-bildern}

\subsubsection{Ein kleines Pendel statt eines unübersichtlichen
Zustandsbergs}\label{ein-kleines-pendel-statt-eines-unuxfcbersichtlichen-zustandsbergs}

Im ursprünglichen Modell gibt es zwei vollständig symmetrische Zustände,
die ausschließlich ineinander umgewandelt werden. Alle anderen denkbaren
Ausgänge heben sich exakt auf. Der ganze Vorgang passt deshalb in eine
Zweizustandsrechnung.

Das ist echte Vereinfachung. Sie erklärt aber noch nicht, wie die
Maschine aus dem Vakuum in einen dieser Startzustände kommt. Dafür fehlt
ein ursprünglicher Vorbereitungsschritt.

\subsubsection{Die richtigen Fenster zeigen mehr als immer mehr Fotos
desselben
Fensters}\label{die-richtigen-fenster-zeigen-mehr-als-immer-mehr-fotos-desselben-fensters}

Einzelne Besetzungsanzeigen sehen nur einen kleinen Teil der inneren
Information. Auch beliebig viele einfache Einzelteilchenmessungen lassen
viel unsichtbar. Geeignet gedrehte gemeinsame Paarmessungen reichen
dagegen für eine vollständige innere Zustandsbeschreibung.

Für die Entwicklung zwischen Paar und Boson brauchen wir zusätzlich
beide Fenster samt ihren tatsächlichen Wahrscheinlichkeiten. Nur das
Paarfenster länger zu beobachten ersetzt die fehlende zweite Ansicht
nicht. Das ist jetzt kein Bauchgefühl mehr, sondern ein genau bestimmter
unsichtbarer Informationsanteil.

\subsubsection{Ein gezielter Anstoß verändert wirklich einen anderen
Ausgang}\label{ein-gezielter-anstouxdf-veruxe4ndert-wirklich-einen-anderen-ausgang}

Wir lassen ein bestimmtes Fermionpaar kurz unter den ursprünglichen
Regeln laufen. Dann ändern wir an einer einzigen Mode nur das Vorzeichen
ihrer Amplitude. Nach weiterer Entwicklung hat eine andere Mode eine
andere Besetzungswahrscheinlichkeit als ohne den Eingriff.

Der Anstoß verändert den anderen Ausgang nicht sofort. Der Unterschied
entsteht erst durch die nachfolgende ursprüngliche Wechselwirkung. Wir
wählen keine günstigen Messergebnisse nachträglich aus und haben keinen
neuen Verbindungsterm in die freie Dynamik eingebaut.

Das ist ein kleiner echter kausaler Effekt im Modell. Die Moden sind
dabei noch keine nachgewiesenen Orte im Raum. Und wir setzen voraus,
dass wir den Anfang, den gezielten Anstoß und die Auslesung tatsächlich
beherrschen.

\subsubsection{Die beiden Aufgaben passen sogar in denselben
Dreierbaustein}\label{die-beiden-aufgaben-passen-sogar-in-denselben-dreierbaustein}

Für genau diesen Versuch reichen drei Zustände: das Anfangspaar, eine
gemeinsame Schwingung der sieben anderen Paare und ein Boson. Wir können
auf diesem selben kleinen Baustein sowohl seinen Zustand auslesen als
auch die Wirkung des Anstoßes berechnen. Das ist die direkteste neue
Verbindung zwischen Schatten und Dynamik.

Die abschließende genaue Einzelmessung kann diese gemeinsame Schwingung
allerdings aufbrechen. Wer danach mit demselben Exemplar weitermachen
will, muss den größeren Zustandsraum wieder mitnehmen oder ein wirklich
anderes, gröberes Messinstrument konstruieren. Nur die angezeigte Zahl
zu kennen heißt noch nicht zu wissen, was die Messung mit der Maschine
gemacht hat.

\subsection{Was die Schattenidee leistet - und was
nicht}\label{was-die-schattenidee-leistet---und-was-nicht}

Deine Idee trifft hier einen prüfbaren Punkt: Verschiedene Ansichten
können gemeinsam verraten, was im Hintergrund passiert. Dafür müssen wir
wissen, wie jede Ansicht aus derselben Maschine entsteht. Dann kann die
eine Ansicht eine fehlende Information der anderen ergänzen.

Das nennen wir hier Rekonstruktion. Noch nicht bewiesen sind eine
räumliche Grenze, eine holografische Beschreibung unserer Welt,
Gravitation oder das vollständige Standardmodell. Ein geometrisches Bild
ist dafür zu wenig; dieselbe Maschine muss auch die richtigen Felder und
Bewegungen erzeugen.

\subsection{Was die letzten Nachträge noch verändert
haben}\label{was-die-letzten-nachtruxe4ge-noch-veruxe4ndert-haben}

Ein zusätzlicher exakt geprüfter Rechenbaustein erlaubt eine stärkere
Untergrenze für die Energie der bekannten niedrigen Entnahmeanregung.
Ihre genaue Energie kennen wir weiter nicht. Wichtig war dabei das
Sortieren: Einige Zahlen im zugeschickten Text gehören zu einem
kleineren Näherungszustand, andere zu einem größeren. Zusammengeschoben
ergeben sie kein einziges berechnetes Teilchen.

Beim Feldtyp gibt es außerdem einen guten Warnhinweis: Eine Aufteilung
kann richtig sein, solange man das System nur dreht, und falsch werden,
wenn Beobachter sich relativ zueinander bewegen. Unser neuer Test zeigt
genau diesen Unterschied für eine bestimmte Lesart des zusammengesetzten
Spinors. Der mathematische Projektor stimmt; ob er zum tatsächlich
gemeinten relativistischen Feld passt, muss der Adapter zeigen.

Der letzte Grundsatztext trägt in seiner Richtung: Die große nächste
Aufgabe ist eine einfache Regel, aus der mehrere unterscheidbare
Vorkommen, ihre gemeinsame Bewegung und ihr Zustand entstehen. Noch mehr
Genauigkeit innerhalb einer einzelnen Bank ersetzt diese Regel nicht.
Eine solche globale Zustandsregel braucht auch keinen äußeren
Experimentator, der erst das ganze Universum vorbereitet.

\subsection{Der jetzt wichtigste nächste
Schritt}\label{der-jetzt-wichtigste-nuxe4chste-schritt}

Nicht noch mehr mögliche Bauteile sammeln, sondern einen dieser kleinen
Versuche vollständig aus dem Compiler herleiten:

\begin{enumerate}
\def\labelenumi{\arabic{enumi}.}
\tightlist
\item
  \textbf{Start:} Woher kommt genau der benötigte Zustand?
\item
  \textbf{Eingriff:} Welche ursprüngliche Operation erzeugt den
  gezielten Handgriff?
\item
  \textbf{Aufzeichnung:} Wie wird der Ausgang gelesen, ohne wichtige
  innere Zusammenhänge unbemerkt zu zerstören?
\end{enumerate}

Danach muss dieselbe Konstruktion auf die bereits abgesicherte geladene
Grundzustandsanregung übertragen werden. Erst wenn diese Anregung unter
kontrollierten Bedingungen zwischen tatsächlich bestimmten Teilen wirkt,
trägt die nächste Frage nach Raum, relativistischen Feldern und
Gravitation.

Die kurze Antwort lautet: Wir haben jetzt kleinere, zusammenhängende
Rechenbausteine und einen echten Wirkungsnachweis. Der noch fehlende
Schlüssel ist ihre ursprüngliche Ausführbarkeit und anschließend der
gemeinsame physikalische Grenzübergang. Das ist konkreter geworden, aber
nicht bereits gelöst.

\end{document}
